\chapter{Iteración 1}

\section{Lista de funcionalidades}
La primera iteración se concibió como la primera entrega jugable del proyecto. Siguiendo
el enfoque de \textit{Feature Driven Development} (FDD), en esta iteración se abordaron
tres funcionalidades que definen el núcleo inicial del sistema:

\begin{enumerate}
\item \textbf{Generación procedural.} Construcción automática de la mazmorra, incluyendo
la topología del nivel, la asignación de salas especiales y el soporte para habitaciones
de distinto tamaño y forma.
\item \textbf{Añadir jugador.} Incorporación de una entidad controlable capaz de aparecer
en la habitación inicial, recorrer el nivel y servir como referencia de navegación y
seguimiento de cámara.
\item \textbf{Añadir enemigos.} Integración de enemigos básicos dentro de las salas, con
comportamiento autónomo simple y con capacidad de condicionar temporalmente la progresión
del jugador mediante el cierre de puertas.
\end{enumerate}

La combinación de estas tres funcionalidades permitió pasar de una base técnica aislada a
una primera versión del bucle jugable: generar un nivel, recorrerlo con el jugador y
resolver encuentros en salas con enemigos.

\section{Diseño}
Desde el punto de vista del diseño, la iteración se estructuró en torno a tres bloques
coordinados: la generación abstracta del mapa, la materialización de ese mapa en salas
jugables y la interacción de entidades dinámicas dentro de dichas salas.

La generación procedural se apoyó en una rejilla lógica para expresar reglas de vecindad,
conectividad y validación espacial. Sobre esa base se definieron salas de distinta forma y
tamaño, evitando limitar el sistema a habitaciones triviales desde el inicio.

Además, se separó la topología del nivel de su representación visual mediante gestores
específicos, de forma que la generación del mapa y la instanciación de salas pudieran
evolucionar con menor acoplamiento.

\section{Implementación}
La implementación del generador procedural partió de una estructura de ocupación del mapa y
de una estrategia de expansión controlada hacia celdas adyacentes. Una vez generada la
topología, se asignaron las salas especiales y se instanciaron las habitaciones jugables
correspondientes.

En paralelo, se integró un controlador básico para el jugador apoyado en
\texttt{Rigidbody2D}, suficiente para validar la escala de movimiento y la navegación por
las salas generadas. También se añadieron enemigos con una máquina de estados elemental y
se implementó el bloqueo y desbloqueo de puertas como mecanismo para encapsular los
encuentros de combate dentro de cada sala.

\section{Resultados}
El resultado de la iteración fue la obtención de una primera base jugable y coherente del
proyecto. Al finalizar esta iteración ya se había conseguido generar un nivel
proceduralmente, recorrerlo con el jugador, poblar salas con enemigos y condicionar el
avance mediante encuentros cerrados.

Como cierre de esta primera entrega, el sistema ya ofrecía un bucle de juego mínimo sobre
el que sostener las iteraciones posteriores.
